% !TEX root = ../../../main.tex
\section{Nondegeneracy}
\label{sec:linear-algebra-i:inner-product-nondegeneracy}

% Source: October 22 handwritten PDF, p. 5.
An inner product allows us to recognize a vector by its pairings with all
other vectors.

\begin{proposition}
  \label{thm:linear-algebra-i:inner-product-nondegeneracy}
  For \(v\in V\),
  \[
    (v,w)=0\text{ for every }w\in V
    \quad\Longleftrightarrow\quad v=0.
  \]
\end{proposition}

\begin{proof}
  If every pairing is zero, take \(w=v\). Then \((v,v)=0\), so
  positive definiteness gives \(v=0\). Conversely, bilinearity gives
  \((0,w)=0\) for every \(w\).
\end{proof}

\begin{proposition}
  \label{thm:linear-algebra-i:inner-product-determines-vector}
  If \(v_1,v_2\in V\) satisfy
  \[
    (v_1,w)=(v_2,w)\qquad\text{for every }w\in V,
  \]
  then \(v_1=v_2\).
\end{proposition}

\begin{proof}
  Bilinearity gives \((v_1-v_2,w)=0\) for every \(w\in V\).
  By \autoref{thm:linear-algebra-i:inner-product-nondegeneracy},
  \(v_1-v_2=0\), hence \(v_1=v_2\).
\end{proof}

\begin{corollary}
  Let \(W\) be a real inner-product space, and let \(L_1,L_2:V\to W\)
  be maps. If
  \[
    (L_1(v),w)=(L_2(v),w)\qquad
    \text{for every }v\in V\text{ and }w\in W,
  \]
  then \(L_1=L_2\).
\end{corollary}
% This consequence is stated on p. 5 without a separate proof.

Such comparisons are useful in the calculus of variations, where an
identity can be tested by pairing with arbitrary vectors or functions.
